% Flat links, curved links, and the geometry of the link field.
% Drafted by Claude Sonnet 5.5 at the author's request, 2026-10-02. Numerical checks: link_curvature_checks/.
\documentclass[11pt]{article}
\usepackage[margin=1in]{geometry}
\usepackage{amsmath,amssymb,amsthm}
\usepackage{booktabs}
\usepackage{array}
\usepackage[T1]{fontenc}
\usepackage{lmodern}
\usepackage{microtype}
\usepackage{xcolor}
\usepackage[colorlinks=true,linkcolor=blue!50!black,citecolor=green!40!black,urlcolor=blue!50!black]{hyperref}
\emergencystretch=2em

\newtheorem{proposition}{Proposition}
\newtheorem{lemma}[proposition]{Lemma}
\newtheorem{corollary}[proposition]{Corollary}
\theoremstyle{definition}
\newtheorem{definition}[proposition]{Definition}
\theoremstyle{remark}
\newtheorem{remark}[proposition]{Remark}

\newcommand{\GL}{\mathrm{GL}}
\newcommand{\gl}{\mathfrak{gl}}
\newcommand{\tr}{\operatorname{tr}}
\newcommand{\im}{\operatorname{im}}
\newcommand{\ad}{\operatorname{ad}}
\newcommand{\cfg}[1]{\texttt{#1}}

\title{Flat Links, Curved Links, and the Geometry of the Link Field:\\Edge-Dependent Generators and Chart-Natural Optimization}
\author{Robert C. Dennis\\\small Technical note, drafted with Claude Sonnet 5.5, October 2, 2026}
\date{}

\begin{document}
\maketitle

\begin{abstract}
A companion note showed that the matrix links of the projected attention routes are group-valued links written in an affine chart, and that with a constant generator the exponential construction adds nothing to the current model. This note asks what an edge-dependent generator can add, and what geometry the link field should be optimized in. We define the triangle holonomy $F_{ijk}=\Theta_{ik}^{-1}\Theta_{ij}\Theta_{jk}$ of a causal link field, whose conjugacy class is gauge invariant. A link field is flat exactly when it is recurrent, that is, when every link is the ordered product of the adjacent links between its endpoints, and for invertible links flat is exactly pure gauge, so the flat transport of any frames is flat, a rotary link is flat, and the frames drop out of the conjugacy class of the holonomy, which is a property of the reference-coordinate field alone. The constant exponential link is maximally curved: its triangle holonomy is the link itself. The routed exponential link has the additive defect of the routing weights as its curvature, and the head-sum link of the model has their multiplicative defect; we show that normalized head-sum routing is flat on strict triangles only for the product form $\beta_{ij}=c_i/c_j$. For a principled edge-dependent link we take the regression of one token's reference state on another's in a joint Gaussian law. The curvature of the regression links is the partial cross-covariance given the middle token, so they are flat exactly when the triple is conditionally independent, as in a Gaussian Markov chain, and their optimal rank-$d$ truncation is reduced-rank regression, which gives the low-rank channel a second-order-statistics origin. Finally, the Fisher metric of a group element acting on a Gaussian belief is the metric already derived for the frame tables; the singular boundary of $\GL(K)$ is at infinite Fisher distance, with the divergence of the Kullback-Leibler divergence growing linearly in $-\log\sigma_{\min}$, whereas additive Euclidean updates in the affine chart cross it freely, and multiplicative group updates never do. All identities were checked in float64 at small dimension; no model was trained, and no claim is made about performance.
\end{abstract}

\section{Setting}
\label{sec:setting}

Tokens sit at positions $0,\dots,N-1$. Under causal support, the source $j$ of a message to token $i$ satisfies $j\le i$, and an edge $(i,j)$ with $j<i$ carries a link $\Theta_{ij}\in\GL(K)$ that transforms as $\Theta_{ij}\mapsto g_i\Theta_{ij}g_j^{-1}$ under a local frame change $g\in\GL(K)^N$. In reference coordinates $\Theta_{ij}=U_iG_{ij}U_j^{-1}$ with $G_{ij}$ invariant, the frames $U_i$ carrying the whole passive action. A companion note \cite{dennis2026linkorigin} showed that the matrix links $L_{ij}=U_iO\,T\,U_j^{-1}$ of the projected routes are, for a channel of this form, the group-valued link $\Omega_{ij}+L_{ij}$ in an affine chart of an embedded $\GL(d)$, and that with a constant generator $X$ the exponential link $U_ie^{X}U_j^{-1}$ is the same operator in logarithmic coordinates. Everything there is edge independent in the reference frame. This note treats what changes when $G_{ij}$ depends on the edge.

\begin{definition}[Triangle holonomy]
For $k<j<i$ let
\begin{equation}
F_{ijk}=\Theta_{ik}^{-1}\Theta_{ij}\Theta_{jk}=U_k\,G_{ik}^{-1}G_{ij}G_{jk}\,U_k^{-1}.
\label{eq:holonomy}
\end{equation}
The field is \emph{flat} when $F_{ijk}=I$ for all triples.
\end{definition}

Under a frame change $F_{ijk}\mapsto g_kF_{ijk}g_k^{-1}$, so the conjugacy class of $F_{ijk}$, and hence its trace, determinant and spectrum, is gauge invariant. $F_{ijk}$ is the holonomy of the closed loop $k\to j\to i\to k$ based at $k$, in which the last step traverses the long edge backward and uses its inverse; this is the sense in which invertibility of the links is what gives a causal field nontrivial loop observables.

\section{Flatness is recurrence}
\label{sec:flat}

\begin{proposition}[Flat is recurrent]
\label{prop:recurrent}
Let $\Theta_{ij}$, $j<i$, be arbitrary $K\times K$ matrices, and call the field flat if $\Theta_{ij}\Theta_{jk}=\Theta_{ik}$ for all $k<j<i$. Then
\begin{equation}
\Theta_{ij}=\Theta_{i,i-1}\Theta_{i-1,i-2}\cdots\Theta_{j+1,j},
\label{eq:recurrent}
\end{equation}
and conversely every choice of adjacent matrices $A_i=\Theta_{i,i-1}$ defines a flat field through the products~\eqref{eq:recurrent}.
\end{proposition}

\begin{proof}
For $i-j\ge2$, flatness with the middle index $i-1$ gives $\Theta_{ij}=\Theta_{i,i-1}\Theta_{i-1,j}$, and induction on $i-j$ gives~\eqref{eq:recurrent}. Conversely, products of adjacent matrices are associative, so $\Theta_{ij}\Theta_{jk}=\Theta_{ik}$.
\end{proof}

\begin{proposition}[Invertible flat is pure gauge]
\label{prop:puregauge}
If every $\Theta_{ij}\in\GL(K)$, the field is flat if and only if there are $g_i\in\GL(K)$ with $\Theta_{ij}=g_ig_j^{-1}$, and then $g_i=\Theta_{i0}$ up to a common right factor. For singular matrices the equivalence fails: the shared projector $\Theta_{ij}=P$ with $P^2=P$ is flat and cannot be written $g_ig_j^{-1}$.
\end{proposition}

\begin{proof}
Set $g_0=I$ and $g_i=\Theta_{i0}=A_i\cdots A_1$. Then $g_ig_j^{-1}=A_i\cdots A_{j+1}=\Theta_{ij}$ by Proposition~\ref{prop:recurrent}, and the converse is the cocycle $g_ig_j^{-1}g_jg_k^{-1}=g_ig_k^{-1}$. The projector satisfies $P\cdot P=P$ but is singular.
\end{proof}

Three consequences follow. First, the flat transport $\Omega_{ij}=U_iU_j^{-1}$ built from arbitrary frames, including frames of the form $U_i=\exp(X_{x_i})\exp(Y_i)$ with positional factors, is pure gauge, and so is any reference-coordinate field of the form $G_{ij}=h_ih_j^{-1}$, for example the rotary link $e^{(p_i-p_j)X}$, because it is absorbed into the frames. Second, the frames drop out of the conjugacy class of the holonomy: by Eq.~\eqref{eq:holonomy}, $F_{ijk}$ is conjugate to $G_{ik}^{-1}G_{ij}G_{jk}$, so curvature is a property of the reference-coordinate field $G_{ij}$ and no choice of frames, positional or not, can flatten it or create it. A constant $G$ is flat only if $G=I$, whatever the positional dependence of the frames. Third, all content of an edge-dependent link beyond a coboundary is curvature, and the conjugacy class of $F_{ijk}$ is the invariant that measures it. This answers the objection, raised in review of the attention manuscript \cite{glkreview2026}, that causal support leaves no invariant loop content: for invertible links the triangle holonomy~\eqref{eq:holonomy} is such content.

\begin{remark}[Recurrent models]
Proposition~\ref{prop:recurrent} says that a flat field is generated by its nearest-neighbor links through path composition. The transition products of a linear state-space recurrence $h_i=A_ih_{i-1}+\dots$, including recurrences whose transitions $A_i$ depend on the input, have exactly this form \cite{gu2022s4,gu2023mamba}, so the recurrent models are the flat link fields and curvature is the deviation of a link between non-adjacent tokens from the composition of the intermediate ones. This is a statement about the cocycle structure of the transition products and not about the models' performance.
\end{remark}

\section{Curvature of the candidate generators}
\label{sec:candidates}

\begin{lemma}[Second-order holonomy]
\label{lem:bch}
For $\Theta_{ij}=U_ie^{X_{ij}}U_j^{-1}$, write $A=X_{ij}$, $B=X_{jk}$, $Z=X_{ik}$. Then $F_{ijk}=U_ke^{-Z}e^{A}e^{B}U_k^{-1}$ and
\begin{equation}
\log\big(e^{-Z}e^{A}e^{B}\big)=(A+B-Z)+\tfrac12\big([A,B]-[Z,A]-[Z,B]\big)+O(3).
\label{eq:bch}
\end{equation}
The first-order term $A+B-Z$ is the discrete exterior derivative $(dX)_{ijk}$ of the algebra-valued $1$-cochain $X$, which vanishes exactly when $X_{ij}=\phi_i-\phi_j$ is exact.
\end{lemma}

\begin{proof}
Apply the Baker-Campbell-Hausdorff series \cite{hall2015} twice: $\log(e^{A}e^{B})=A+B+\tfrac12[A,B]+\dots$, and then $\log(e^{-Z}e^{C})=C-Z-\tfrac12[Z,C]+\dots$ with $C=A+B+\tfrac12[A,B]$.
\end{proof}

\begin{table}[ht]
\centering
\small
\begin{tabular}{>{\raggedright}p{0.30\linewidth}>{\raggedright}p{0.36\linewidth}>{\raggedright\arraybackslash}p{0.26\linewidth}}
\toprule
Link field & Holonomy $F_{ijk}$ & Flat when \\
\midrule
Positional, $U_ie^{(p_i-p_j)X}U_j^{-1}$ & $I$ & always \\
Constant, $U_ie^{X}U_j^{-1}$ & $U_ke^{X}U_k^{-1}$ & only if $e^{X}=I$ \\
Routed exponent, $X_{ij}=OB_{ij}T$, $TO=I$ & $U_kO\,e^{B_{ij}+B_{jk}-B_{ik}}T\,U_k^{-1}$ & $B$ additive: $B_{ij}+B_{jk}=B_{ik}$ \\
Head sum, $U_iOB_{ij}TU_j^{-1}$, $TO=I$ & $U_kO\,\frac{B_{ij}B_{jk}}{B_{ik}}\,T\,U_k^{-1}$ & $B$ multiplicative: $B_{ij}B_{jk}=B_{ik}$ \\
\bottomrule
\end{tabular}
\caption{Triangle holonomy of the candidate edge-dependent link fields. The routing matrices $B_{ij}=\bigoplus_h\beta_{ij}^{(h)}I_d$ act headwise.}
\label{tab:candidates}
\end{table}

Table~\ref{tab:candidates} collects the exact holonomies. The first row is Proposition~\ref{prop:puregauge}. The second shows that the constant link of the companion note is the most curved edge-independent choice: every triangle holonomy equals the link itself, whose first-order part is $X=(dX)_{ijk}$, which is why $L$ carries curvature and is not gauge. For the routed exponent and the head sum with $TO=I$, the head blocks commute and the formulas follow by multiplying block by block; in the routed exponent the generators add, in the head sum the weights multiply.

For the routing $\beta_{ij}^{(h)}=\pi_{ij}^{(h)}e^{-D_{ij}^{(h)}/\kappa}/Z_i^{(h)}$, the logarithm of the multiplicative defect is
\begin{equation}
\log\frac{\beta_{ij}\beta_{jk}}{\beta_{ik}}=-\frac{D_{ij}+D_{jk}-D_{ik}}{\kappa}+\log\frac{\pi_{ij}\pi_{jk}}{\pi_{ik}}-\log Z_j,
\label{eq:kl_defect}
\end{equation}
as in the companion note, where $Z_i$ cancels between $\beta_{ij}$ and $\beta_{ik}$.

\begin{proposition}[Flat normalized routing]
\label{prop:flatrouting}
Let $\beta_{ij}>0$ for $j\le i$ with $\sum_{j\le i}\beta_{ij}=1$. The head-sum link with $TO=I$ is flat on strict triangles $k<j<i$ if and only if $\beta_{ij}=c_i/c_j$ for $0<j<i$ and $\beta_{i0}=c_i$, where $c_i:=\beta_{i0}$ and $c_0=\beta_{00}=1$.
\end{proposition}

\begin{proof}
By Table~\ref{tab:candidates}, flatness is $\beta_{ij}\beta_{jk}=\beta_{ik}$ for each head and strict triple. Taking $k=0$ gives $\beta_{ij}\beta_{j0}=\beta_{i0}$, so $\beta_{ij}=c_i/c_j$ with $c_i=\beta_{i0}$, and $c_0=\beta_{00}=1$ because the first row has the single source $0$. Conversely this form satisfies $\beta_{ij}\beta_{jk}=c_ic_j/(c_jc_k)=\beta_{ik}$.
\end{proof}

The form $\beta_{ij}=c_i\,w_j$ with $w_j=1/c_j$ is a separable attention pattern in which the weight on source $j$ is the reciprocal of $j$'s own first-source weight. A valid example exists, for instance $c=(1,\,0.5,\,0.2,\,0.05,\,0.01)$ with the diagonal weights absorbing the remaining mass, but it is a thin subset: uniform causal attention violates the cocycle already on five tokens (largest violation $1/6$), and generic content-dependent routing has a nonzero multiplicative defect. This is a statement about the head-sum link of Option C of the companion note, not about attention in general.

\section{Regression links: an inference-derived edge-dependent link}
\label{sec:regression}

The generators above are postulated. A link that is derived from a probabilistic model is the best linear predictor of one token's reference state from another's. Let $(\xi_0,\dots,\xi_{N-1})$ be jointly Gaussian in reference coordinates, centered for simplicity (nonzero means add an intercept), with $\Sigma_{ab}=\mathrm{Cov}(\xi_a,\xi_b)$ and invertible marginals $\Sigma_{bb}$, and define
\begin{equation}
W_{ab}=\Sigma_{ab}\Sigma_{bb}^{-1}=\arg\min_{W}\mathbb E\,\|\xi_a-W\xi_b\|^2_{\Sigma_{aa}^{-1}}\ \text{over unrestricted }W.
\label{eq:regression}
\end{equation}
Under a frame change $\xi_a\mapsto g_a\xi_a$ one has $\Sigma_{ab}\mapsto g_a\Sigma_{ab}g_b^\top$ and $W_{ab}\mapsto g_aW_{ab}g_b^{-1}$, so $\Theta_{ab}=W_{ab}$ is a link with the correct covariance law.

\begin{proposition}[Curvature is partial covariance]
\label{prop:markov}
For $k<j<i$ let $\Sigma_{ik|j}=\Sigma_{ik}-\Sigma_{ij}\Sigma_{jj}^{-1}\Sigma_{jk}$. Then
\begin{equation}
W_{ij}W_{jk}=W_{ik}-\Sigma_{ik|j}\Sigma_{kk}^{-1},\qquad F_{ijk}=I-W_{ik}^{-1}\Sigma_{ik|j}\Sigma_{kk}^{-1}
\label{eq:partialcov}
\end{equation}
when $W_{ik}$ is invertible. Hence $F_{ijk}=I$ if and only if $\Sigma_{ik|j}=0$, which for a Gaussian law is the conditional independence $\xi_i\perp\xi_k\mid\xi_j$ \cite{lauritzen1996}. In particular, if the sequence is a Gaussian vector Markov chain $\xi_i=A\xi_{i-1}+\varepsilon_i$ with independent noise, then the regression links are flat and $W_{ij}=A^{\,i-j}$.
\end{proposition}

\begin{proof}
$W_{ij}W_{jk}=\Sigma_{ij}\Sigma_{jj}^{-1}\Sigma_{jk}\Sigma_{kk}^{-1}=(\Sigma_{ik}-\Sigma_{ik|j})\Sigma_{kk}^{-1}$, which is the first identity, and $F_{ijk}=W_{ik}^{-1}W_{ij}W_{jk}$ gives the second. For the chain, $\Sigma_{ij}=A^{\,i-j}\Sigma_{jj}$ for $j\le i$, so $W_{ij}=A^{\,i-j}$ and $W_{ij}W_{jk}=W_{ik}$.
\end{proof}

The curvature of the regression links is the information that token $i$ carries about token $k$ beyond what passes through token $j$. In this reading a flat link field is a transmission chain in which each token depends on the earlier ones only through its predecessor, and curvature is the amount of direct, unmediated dependence, the thing attention exists to provide. The converse of the Markov statement is not claimed: vanishing of the partial covariance for every triple with a single intermediate token is weaker than the full Markov property.

\paragraph{The channel as reduced-rank regression.}
The rank-$d$ structure of the model's links has an inference-derived counterpart. Whiten both sides and let $\Sigma_{ii}^{-1/2}\Sigma_{ij}\Sigma_{jj}^{-1/2}=\sum_sa_su_sv_s^\top$ be a singular value decomposition with $a_1\ge a_2\ge\cdots$, the canonical correlations.

\begin{proposition}[Reduced-rank regression]
\label{prop:rrr}
Among matrices of rank at most $d$, the minimizer of $\mathbb E\|\xi_i-W\xi_j\|^2_{\Sigma_{ii}^{-1}}$ is
\begin{equation}
W_d=\Sigma_{ii}^{1/2}\Big(\sum_{s\le d}a_su_sv_s^\top\Big)\Sigma_{jj}^{-1/2},
\label{eq:rrr}
\end{equation}
and its excess risk over the full regression $W_{ij}$ is $\sum_{s>d}a_s^2$.
\end{proposition}

\begin{proof}
The risk equals $\|\Sigma_{ii}^{-1/2}(W-W_{ij})\Sigma_{jj}^{1/2}\|_F^2$ plus the risk of $W_{ij}$, and the Eckart-Young theorem gives the nearest rank-$d$ matrix in the whitened Frobenius norm \cite{izenman1975,reinsel1998}.
\end{proof}

A factorization $W_d=OT$ with $O$ of size $K\times d$ and $T$ of size $d\times K$ is therefore the optimal communicating channel of width $d$, with the leading canonical directions as $\operatorname{im}O$ and the row space of $T$ as the sending directions. With $K=Hd$ the heads would partition the canonical directions into $H$ blocks of $d$. This is an interpretation: the model's heads are not trained to this criterion, and the identification of $(O,T)$ with the reduced-rank solution is a hypothesis to test, not a result.

\begin{remark}[What the regression link needs]
Equation~\eqref{eq:regression} requires the cross-covariances $\Sigma_{ab}$, which vanish in the mean-field family $q=\prod_iq_i$ of the model, where the regression link is zero. A correlated variational family, for example a Gaussian with sparse block precision, or a generative model whose prior couples the tokens, would supply them; the learned links $O\,T$ of the projected routes can be read as an amortized replacement.
\end{remark}

\section{Chart-natural optimization and the singular boundary}
\label{sec:natural}

A group element $g\in\GL(K)$ acts on a Gaussian belief by $N(\mu,\Sigma)\mapsto N(g\mu,g\Sigma g^\top)$. Write a tangent vector as $\delta g=Ag$ with $A=\delta g\,g^{-1}\in\gl(K)$, and let $(\mu',\Sigma')=(g\mu,g\Sigma g^\top)$ be the output with Cholesky factor $L'$.

\begin{proposition}[Fisher metric of the link]
\label{prop:fisher}
The Fisher-Rao metric pulled back through the action is
\begin{equation}
ds^2=\|Z\eta\|^2+2\,\|\tfrac12(Z+Z^\top)\|_F^2,\qquad Z=L'^{-1}AL',\ \eta=L'^{-1}\mu',
\label{eq:fisher}
\end{equation}
which depends on the output belief only. In the exponential chart $g=e^{X}$ one has $A=\int_0^1e^{s\,\ad_X}\delta X\,ds=\frac{e^{\ad_X}-1}{\ad_X}\delta X$.
\end{proposition}

\begin{proof}
The tangent of the output is $\delta\mu'=A\mu'$, $\delta\Sigma'=A\Sigma'+\Sigma'A^\top$. The Fisher metric of a Gaussian is $\delta\mu'^\top\Sigma'^{-1}\delta\mu'+\tfrac12\tr(\Sigma'^{-1}\delta\Sigma'\Sigma'^{-1}\delta\Sigma')$; substituting $\Sigma'=L'L'^\top$ gives~\eqref{eq:fisher}. The derivative of the exponential map is standard \cite{hall2015}.
\end{proof}

Equation~\eqref{eq:fisher} is the metric derived for the frame tables in the companion note on frame-chart geometry \cite{dennis2026phinatgrad}, now for a link acting on a source belief. For the model, the relevant metric on an edge averages~\eqref{eq:fisher} over the source beliefs with the routing weights.

\begin{proposition}[Infinite distance to the singular boundary]
\label{prop:barrier}
Let $g_t=\mathrm{diag}(e^{-t},1,\dots,1)$ act on $N(\mu,I)$, so $\sigma_{\min}(g_t)=e^{-t}$. Then
\begin{equation}
\mathrm{KL}\big(q_t\,\|\,q_0\big)=t+\tfrac12\big(e^{-2t}-1+(1-e^{-t})^2\mu_1^2\big)\sim t,
\label{eq:klbarrier}
\end{equation}
the path has Fisher length $\sqrt{\mu_1^2+2}\,t$, and the Fisher-Rao distance from $q_0$ to $q_t$ is at least $\sqrt2\,t$.
\end{proposition}

\begin{proof}
The Kullback-Leibler divergence follows from the Gaussian formula with $\Sigma_t=\mathrm{diag}(e^{-2t},1,\dots)$ and $\mu_t=g_t\mu$. Along the path $A=-E_{11}$ and $L'=g_t$, so $Z=-E_{11}$ and $\eta=\mu$, which gives $ds^2=(\mu_1^2+2)dt^2$ from~\eqref{eq:fisher}. The lower bound follows from the univariate Fisher metric $(dm^2+2\,d s^2)/s^2$, which gives $d\ge\sqrt2\,|\log(s_1/s_0)|$ for any means, combined with monotonicity of the Fisher-Rao distance under marginalization onto the first coordinate \cite{amari2000}.
\end{proof}

The singular boundary of $\GL(K)$ is therefore at infinite Fisher distance from every interior point, and the divergence grows linearly in $-\log\sigma_{\min}$. In an affine chart with Euclidean metric the same boundary is at finite distance $\sigma_{\min}$.

\begin{corollary}[Updates and the boundary]
\label{cor:updates}
An additive update $M\leftarrow M-\eta\Delta$ in the affine chart can cross $\det=0$ in a single step, by jumping to the other side of the boundary, whenever $\eta\|\Delta\|_2\ge\sigma_{\min}(M)$ and $\Delta$ is suitably aligned with the smallest singular direction; a multiplicative update $g\leftarrow e^{-\eta A}g$ never does, because $\det(e^{-\eta A}g)=e^{-\eta\tr A}\det g\ne0$. A natural-gradient step, whose Fisher length is proportional to $\eta$, needs at least $\sqrt2\,|\log\sigma_{\min}|$ divided by that length steps to approach the boundary, by Proposition~\ref{prop:barrier}.
\end{corollary}

For the trained run measured in the companion note, the content learning rate is $2.5\times10^{-3}$ per element under AdamW, and the stacked operator has $\sigma_{\min}$ between $10^{-5}$ and $4\times10^{-3}$, below or comparable to the scale of a single step times $\|T\|\sqrt K$, about $3\times10^{-2}$ for $\|T\|\approx1.5$ and $K=64$. This is an order-of-magnitude comparison and is not a measurement of the crossings. The affine-chart updates are consistent with the observed repeated sign changes of $\det M$, whereas group updates would not allow them. None of this bears on the link that run actually applied: the residual-augmented operator $I+M$ stayed far from the boundary. The barrier matters for designs that apply a bare group element, such as a residual-free route with tiled heads or a group-valued link without a flat term.

\paragraph{Where the barrier comes from.}
The loss does not diverge at the boundary, since the predictive and receiver laws carry covariance floors and noise floors, so the barrier is a property of the geometry used to move and not of the objective. The vfe3 code already uses group updates for the frame tables (\cfg{m\_phi\_update\_mode} with \cfg{pullback\_group} and \cfg{fisher\_group}), and the same construction, a multiplicative retraction with a Fisher preconditioner, could be applied to the content maps.

\section{What follows for the model}
\label{sec:model}

Three conclusions follow. First, curvature is the frame-independent content of the link field, and a coboundary is not content. Whatever the learned link does beyond a pure-gauge factor, including the constant $G$ of the current model, is measured by the holonomy $F_{ijk}$, whose trace and determinant can be computed for the head-sum link directly from logged routing weights through~\eqref{eq:kl_defect}. A useful baseline is the flat pattern $\beta_{ij}=c_i/c_j$ of Proposition~\ref{prop:flatrouting} and the uniform pattern, whose defect is nonzero. Second, a principled non-positional edge-dependent link is the regression link of Section~\ref{sec:regression}: it is covariant, its curvature has a clear probabilistic meaning, and its channel-restricted form is reduced-rank regression. Using it requires a correlated variational family. Third, if a design applies a bare group element, the optimizer should either update multiplicatively or move in a metric that sees the boundary; additive affine updates give no protection.

These are claims about the geometry of candidate designs. None has been tested in training, and the relationship of any of them to perplexity is unknown.

\section{Numerical checks}
\label{sec:checks}

The script \texttt{link\_curvature\_checks/checks.py}, with output in \texttt{checks\_output.txt}, evaluates each identity in float64 at $K=4$ with random invertible matrices. Recurrent links satisfy $F_{ijk}=I$ to $1.8\times10^{-15}$ and the pure-gauge form $g_ig_j^{-1}$ to $1.1\times10^{-15}$; the shared rank-one projector satisfies the cocycle exactly and has rank $1$. The holonomy of the constant link equals $U_ke^{X}U_k^{-1}$ to $2.0\times10^{-15}$, the positional link has $\max|F-I|=3.3\times10^{-13}$, the routed exponent matches its additive formula to $2.2\times10^{-15}$, and the head sum matches its multiplicative formula to $2.1\times10^{-15}$. The Baker-Campbell-Hausdorff expansion of Lemma~\ref{lem:bch} has first-order error $3.6\times10^{-2}$, $9.1\times10^{-3}$, $2.3\times10^{-3}$ and second-order error $3.6\times10^{-3}$, $4.6\times10^{-4}$, $5.7\times10^{-5}$ at generator scales $0.1$, $0.05$, $0.025$, scaling as $\varepsilon^2$ and $\varepsilon^3$. The flat routing pattern with $c=(1,0.5,0.2,0.05,0.01)$ has all weights positive, row sums one, and cocycle defect $0$, while uniform causal attention has defect $1/6$. For a generic Gaussian, Eq.~\eqref{eq:partialcov} holds to $1.4\times10^{-16}$ and the second identity to $5.3\times10^{-16}$; a generic joint law has $\max|F-I|=3.4$, while an AR(1) chain has $1.2\times10^{-15}$ and $W_{ij}=A^{\,i-j}$ to $1.1\times10^{-16}$; the covariance law $W_{ab}\mapsto g_aW_{ab}g_b^{-1}$ holds with distinct per-token frames to $1.3\times10^{-15}$. The reduced-rank excess risk $0.096854$ equals the sum of squared discarded singular values $0.096854$, and no factor perturbation or random rank-$2$ matrix does better. The Fisher expansion $2\,\mathrm{KL}/(t^2ds^2)$ equals $1.0054$, $1.0027$, $1.0013$ at $t=10^{-2},5\times10^{-3},2.5\times10^{-3}$, converging to one linearly in $t$; the derivative of the exponential map is verified to $8.9\times10^{-9}$ by finite differences; and along the shrinking path the divergence is $2.133,4.499,6.809,9.112$ at $\sigma_{\min}=10^{-1},\dots,10^{-4}$, with increments per decade of $2.37$, $2.31$, $2.30$ approaching $\ln10=2.303$. These checks ground the algebra at one small size and say nothing about trained models.

\begin{thebibliography}{99}
\bibitem{dennis2026linkorigin} R.~C.~Dennis, \emph{Where the matrix links come from: group-valued links and their communicating parts}, technical note, manuscripts/Attention/link\_origin.tex, 2026.
\bibitem{dennis2026phinatgrad} R.~C.~Dennis, \emph{Frame-chart geometry for the stored frame tables: the Frobenius-pullback step and the belief-Fisher natural gradient}, technical note, manuscripts/Attention/2026-09-28-phi-natgrad.tex, 2026.
\bibitem{glkreview2026} Referee panel report on the GL(K) attention v2 manuscript, manuscripts/reviews/2026-09-23-glk-v2/REVIEW.md, 2026.
\bibitem{amari2000} S.~Amari and H.~Nagaoka, \emph{Methods of Information Geometry}, American Mathematical Society and Oxford University Press, 2000.
\bibitem{gu2022s4} A.~Gu, K.~Goel, and C.~R{\'e}, Efficiently modeling long sequences with structured state spaces, \emph{International Conference on Learning Representations}, 2022.
\bibitem{gu2023mamba} A.~Gu and T.~Dao, Mamba: linear-time sequence modeling with selective state spaces, arXiv:2312.00752, 2023.
\bibitem{hall2015} B.~C.~Hall, \emph{Lie Groups, Lie Algebras, and Representations: An Elementary Introduction}, 2nd ed., Graduate Texts in Mathematics 222, Springer, 2015.
\bibitem{izenman1975} A.~J.~Izenman, Reduced-rank regression for the multivariate linear model, \emph{Journal of Multivariate Analysis} 5, 248--264, 1975.
\bibitem{lauritzen1996} S.~L.~Lauritzen, \emph{Graphical Models}, Oxford University Press, 1996.
\bibitem{reinsel1998} G.~C.~Reinsel and R.~P.~Velu, \emph{Multivariate Reduced-Rank Regression: Theory and Applications}, Lecture Notes in Statistics 136, Springer, 1998.
\end{thebibliography}

\end{document}
